\documentclass[12pt,a4paper]{article}
\usepackage[utf8]{inputenc}
\usepackage[T1]{fontenc}
\usepackage[english]{babel}
\usepackage{amsmath,amssymb,amsthm}
\usepackage{mathtools}
\usepackage{geometry}
\usepackage{titlesec}
\usepackage{fancyhdr}
\usepackage{hyperref}

\geometry{margin=1in}

\pagestyle{fancy}
\fancyhf{}
\fancyhead[L]{\small Cayley, Collected Mathematical Papers, Vol.~IX}
\fancyhead[R]{\small Pages 301--350}
\fancyfoot[C]{\thepage}
\renewcommand{\headrulewidth}{0.4pt}

\newtheoremstyle{cayley}% name
  {\topsep}% Space above
  {\topsep}% Space below
  {\itshape}% Body font
  {}% Indent amount
  {\bfseries}% Theorem head font
  {.}% Punctuation after theorem head
  {.5em}% Space after theorem head
  {}% Theorem head spec

\theoremstyle{cayley}
\newtheorem{art}{Art.}[section]

\titleformat{\section}{\large\bfseries\centering}{\thesection.}{0.5em}{}
\titleformat{\subsection}{\normalsize\bfseries}{\thesubsection.}{0.5em}{}

\DeclareMathOperator{\arccot}{arccot}

\begin{document}

\begin{center}
{\LARGE\bfseries Collected Mathematical Papers\\[0.5em]}
{\Large Arthur Cayley, Vol.~IX\\[1em]}
{\large Pages 301--350\\[0.5em]}
{\normalsize Papers 603--607}
\end{center}

\vspace{2em}

\tableofcontents
\newpage


% ====================================================================
% PAPER 603
% ====================================================================
\newpage
\section[On the Potential of the Ellipse and the Circle]%
{603. On the Potential of the Ellipse and the Circle.}

\begin{center}
{\small [From the \textit{Proceedings of the London Mathematical Society}, vol.~VI. (1874--1875), pp.~38--58.\\
Read January 14, 1875.]}
\end{center}

\subsection*{The Potential of the Ellipse.}

\begin{art}
I consider the potential of an ellipse (or say an elliptic plate of uniform density); viz. this is
\[
V = \iint \frac{dx\,dy}{\sqrt{(a-x)^2 + (b-y)^2 + c^2}},
\]
the limits being given by the equation
\[
\frac{x^2}{f^2} + \frac{y^2}{g^2} = 1.
\]
Writing herein
\[ x = mf\cos u, \quad y = mg\sin u, \]
we have $dx\,dy = fg\,m\,dm\,du$; and consequently
\[
V = \iint \frac{fg\,m\,dm\,du}{\sqrt{(a - mf\cos u)^2 + (b - mg\sin u)^2 + c^2}},
\]
where the integrations are to be taken from $m=0$ to $m=1$, and from $u=0$ to $u=2\pi$.
\end{art}

\begin{art}
It is to be remarked that, by first performing the integration in regard to $m$, we may reduce the potential to the form $\int du\,F$, where $F$ is an algebraic function of $\cos u$, $\sin u$; and that the result so obtained, although in the general case too complex to be manageable, is a useful one in the case $f=g$, where the ellipse becomes a circle. The case of the circle will be treated of separately, but in the general case it will be sufficient to show that the integral is of the form in question.
\end{art}

\begin{art}
To accomplish this, writing
\[ A = a^2 + b^2 + c^2, \quad B = af\cos u + bg\sin u, \quad C = f^2\cos^2 u + g^2\sin^2 u, \]
then the integral in regard to $m$ is
\[
\int \frac{m\,dm}{\sqrt{A - 2Bm + Cm^2}},
\]
which is
\[
= \frac{1}{C}\sqrt{A - 2Bm + Cm^2} + \frac{B}{C^{\frac{3}{2}}}\log(Cm - B + \sqrt{C}\sqrt{A - 2Bm + Cm^2}).
\]
Taken between the limits $0$ and $1$, this is
\[
= \frac{1}{C}(\sqrt{A - 2B + C} - \sqrt{A}) + \frac{B}{C^{\frac{3}{2}}}\log\frac{C - B + \sqrt{C}\sqrt{A - 2B + C}}{-B + \sqrt{AC}};
\]
and we have therefore
\[
V = fg\int du\, \frac{1}{C}(\sqrt{A - 2B + C} - \sqrt{A}) + fg\int du\, \frac{B}{C^{\frac{3}{2}}}\log\frac{C - B + \sqrt{C}\sqrt{A - 2B + C}}{-B + \sqrt{AC}},
\]
where, for greater clearness, the value of the coefficient $\frac{B}{C^{\frac{3}{2}}}$ of the logarithmic term has been written at full length.
\end{art}

\begin{art}
But this coefficient admits of algebraic integration, viz. we have
\[
\int \frac{B\,du}{C^{\frac{3}{2}}}
= \frac{ag\sin u - bf\cos u}{(f^2\cos^2 u + g^2\sin^2 u)^{\frac{1}{2}}};
\]
hence, integrating the second term by parts, we have
\begin{multline*}
V = fg\int du\, \frac{1}{C}(\sqrt{A - 2B + C} - \sqrt{A}) \\
+ \frac{ag\sin u - bf\cos u}{(f^2\cos^2 u + g^2\sin^2 u)^{\frac{1}{2}}}\log T
- \int du\, \frac{ag\sin u - bf\cos u}{(f^2\cos^2 u + g^2\sin^2 u)^{\frac{1}{2}}} \cdot \frac{1}{T}\frac{dT}{du},
\end{multline*}
where the second term, taken between the limits $u=0$, $u=2\pi$, is $=0$; and $\frac{1}{T}\frac{dT}{du}$ being an algebraic function of $\sin u$, $\cos u$, the potential is expressed in the form in question.
\end{art}

\begin{art}
But we may, by means of a transformation upon $u$ (that made use of in Gauss' Memoir\footnote{[Ges.\ Werke, t.~III., pp.~333--355; in particular, l.c., p.~338].} on the attraction of an elliptic ring), transform the expression so as to obtain the integral in regard to $m$ under a much more simple form. We, in fact, assume
\begin{align*}
\cos u &= \frac{\alpha + \alpha'\cos\tau + \alpha''\sin\tau}{\gamma + \gamma'\cos\tau + \gamma''\sin\tau}, \\
\sin u &= \frac{\beta + \beta'\cos\tau + \beta''\sin\tau}{\gamma + \gamma'\cos\tau + \gamma''\sin\tau},
\end{align*}
where the nine coefficients are such that identically
\[
(\alpha + \alpha'\cos\tau + \alpha''\sin\tau)^2 + (\beta + \beta'\cos\tau + \beta''\sin\tau)^2 - (\gamma + \gamma'\cos\tau + \gamma''\sin\tau)^2 = \cos^2\tau + \sin^2\tau - 1,
\]
(this of course renders the two equations consistent); and also that
\[
(a - mf\cos u)^2 + (b - mg\sin u)^2 + c^2 = \frac{G + G'\cos^2\tau + G''\sin^2\tau}{(\gamma + \gamma'\cos\tau + \gamma''\sin\tau)^2}.
\]
This last condition gives, for the determination of the coefficients $G$, $G'$, $G''$, the identity
\[
\frac{a^2}{G + m^2f^2} + \frac{b^2}{G + m^2g^2} + \frac{c^2}{G} = 1,
\]
or, what is the same thing,
\[
\frac{f^2}{G + m^2f^2} + \frac{g^2}{G + m^2g^2} + \frac{c^2}{G} - 1 = 0.
\]
This equation has one positive root, which may be taken to be $G$, and two negative roots, which will then be $-G'$, $-G''$; viz. $G$, $G'$, $G''$ are thus all positive; and $G$ denotes the positive root of the last-mentioned equation.
\end{art}

\begin{art}
We have
\[ du = \frac{d\tau}{(G + G'\cos^2\tau + G''\sin^2\tau)^{\frac{1}{2}}}, \]
and thence
\[
V = fg\int m\,dm\int \frac{d\tau}{(G + G'\cos^2\tau + G''\sin^2\tau)^{\frac{3}{2}}},
\]
the integral in regard to $\tau$ being taken from $0$ to $2\pi$, or, what is the same thing, we may multiply by $4$ and take the integral only from $0$ to $\frac{\pi}{2}$; viz. we thus have
\[
V = 4fg\int m\,dm\int_0^{\frac{\pi}{2}} \frac{d\tau}{(G + G'\cos^2\tau + G''\sin^2\tau)^{\frac{3}{2}}},
\]
where the integral in regard to $\tau$ can be at once reduced to the standard form of an elliptic function, or it might be calculated by Gauss' method of the arithmetico-geometrical mean.
\end{art}

\begin{art}
But, for the present purpose, a further reduction is required. Writing
\[ t = G + (G + G')\cos^2\tau, \]
we have
\begin{align*}
t + G' &= (G + G')\frac{\sin^2\tau}{\sin^2 i}, \\
t + G'' &= (G + G'\cos^2\tau + G''\sin^2\tau)\frac{\sin^2 i}{\sin^2\tau};
\end{align*}
whence
\[
\sqrt{t - G \cdot t + G' \cdot t + G''} = (G + G')(G + G'\cos^2\tau + G''\sin^2\tau)^{\frac{1}{2}}\frac{\sin^2 i}{\sin^2\tau};
\]
moreover
\[ dt = -2(G + G')^{\frac{1}{2}}\frac{\sin^2 i}{\sin^2\tau}\,d\tau. \]
Hence
\[
\frac{dt}{\sqrt{t - G \cdot t + G' \cdot t + G''}} = \frac{-2\,d\tau}{(G + G'\cos^2\tau + G''\sin^2\tau)^{\frac{3}{2}}};
\]
and, observing that to the limits $0$, $\frac{\pi}{2}$ of $\tau$ correspond the limits $\infty$, $G$ of $t$, we thence obtain
\[
V = 2fg\int m\,dm\int_G^{\infty} \frac{dt}{\sqrt{t - G \cdot t + G' \cdot t + G''}};
\]
or, what is the same thing,
\[
V = 2fg\int m\,dm\int_0^{\infty} \frac{dt}{\sqrt{t(t + m^2f^2)(t + m^2g^2)}},
\]
where $G$ denotes, as before, the positive root of the equation
\[
\frac{a^2}{G + m^2f^2} + \frac{b^2}{G + m^2g^2} + \frac{c^2}{G} = 1.
\]
\end{art}

\begin{art}
Writing for $f$, $mf$, and for $G$, $m^2G$, the formula becomes
\[
V = 2fg\int m\,dm\int_G^{\infty} \frac{dt}{\sqrt{t(t + f^2)(t + g^2)}},
\]
where $G$ now denotes the positive root of the equation
\[
\frac{a^2}{G + f^2} + \frac{b^2}{G + g^2} + \frac{c^2}{G} = 1.
\]
\end{art}

\begin{art}
Thus $G$ is a function of $m$; but it is to be remarked that the integration in respect to $m$ can be performed through the integral sign $\int\frac{dt}{\sqrt{}}$ in precisely the same way as if $G$ were constant, and that we, in fact, have
\[
V = 2fg\int_0^1 m\,dm\int_G^{\infty} \frac{dt}{\sqrt{t(t + m^2f^2)(t + m^2g^2)}},
\]
where the function of $m$ is to be taken between the limits $0$ and $1$.

The reason is that, differentiating this last integral in respect to $m$, the term depending on the variation of the limit $G$ is
\[
\sqrt{\frac{1}{G + m^2f^2} + \frac{1}{G + m^2g^2} - \frac{1}{G}} \cdot \frac{dG}{dm},
\]
which is $=0$ in virtue of the equation which defines $G$; hence the whole result is the term arising from the variation of $m$ in so far as it appears explicitly.
\end{art}

\begin{art}
Proceeding next to take the function of $m$ between the two limits: for $m=0$ we have $G=\infty$, and the integral vanishes; for $m=1$ we have $G$ the positive root of the equation
\[
\frac{a^2}{G + f^2} + \frac{b^2}{G + g^2} + \frac{c^2}{G} = 1,
\]
or, using $\theta$ to denote the positive root of this equation, the value is $G = \theta$; we thus finally obtain
\[
V = 2fg\int_{\theta}^{\infty} \frac{dt}{\sqrt{t(t + f^2)(t + g^2)}},
\]
as the expression for the potential of the ellipse semiaxes $(f,g)$ on the point $(a,b,c)$.
\end{art}

\subsection*{Case where the Attracted Point is on the Focal Hyperbola.}

\begin{art}
The result becomes very simple when the attracted point is in the focal hyperbola of the ellipse, viz. when we have $b=0$ and $\frac{a^2}{f^2 - g^2} = 1$. The function
\[
\frac{1}{\sqrt{t(t + f^2)(t + g^2)}}
\]
is here
\[
\frac{1}{\sqrt{t(t + f^2)(t + g^2)}} = \frac{1}{t\sqrt{t + f^2}} - \frac{g^2}{t(t + f^2)\sqrt{t + g^2}}.
\]
Hence also
\[
\frac{d\theta}{\sqrt{\theta(\theta + f^2)(\theta + g^2)}} = \frac{d\theta}{\theta\sqrt{\theta + f^2}} - \frac{g^2\,d\theta}{\theta(\theta + f^2)\sqrt{\theta + g^2}};
\]
introducing this value, the function in question becomes
\[
\int_{\theta}^{\infty} \left(\frac{1}{t\sqrt{t + f^2}} - \frac{g^2}{t(t + f^2)\sqrt{t + g^2}}\right) dt
\]
and we have
\begin{align*}
V &= 2fg\int_{\theta}^{\infty} \frac{dt}{t\sqrt{t + f^2}} - 2fg\int_{\theta}^{\infty} \frac{g^2\,dt}{t(t + f^2)\sqrt{t + g^2}} \\
&= \frac{2fg}{f}\left[\frac{1}{\sqrt{t}}\tan^{-1}\frac{\sqrt{t}}{f} - \frac{g}{f\sqrt{t}}\tan^{-1}\frac{g\sqrt{t}}{f\sqrt{t + g^2}}\right]_{\theta}^{\infty},
\end{align*}
which, writing $t = f^2\tan^2\phi$, becomes
\begin{align*}
V &= \frac{2g}{f}\left(\frac{\pi}{2} - \tan^{-1}\frac{f}{\sqrt{\theta}}\right) \\
&= \frac{2g}{f}\left(\frac{\pi}{2} - \frac{f}{\sqrt{\theta}}\right) \\
&= 2\pi\left(\frac{g}{f}\sqrt{c^2 + f^2} - c\right);
\end{align*}
or, substituting for $\theta$ its value $\frac{f^2c^2}{g^2}$, this is
\[
V = 2\pi(\sqrt{c^2 + g^2} - c),
\]
which is, in fact, the potential of the circle $x^2 + y^2 = g^2$ on the axial point $(0, 0, c)$ and, observing that the value is independent of $f$, we have at once the theorem that, considering $f$ as variable, and taking the attracted point at the constant altitude $c$ in the focal hyperbola $\frac{x^2}{f^2} - \frac{y^2}{g^2} = 1$, the potential is the same, whatever is the value of the semi-axis major $f$ of the ellipse.
\end{art}


\begin{art}
A point in the focal hyperbola determines, with the ellipse, a right circular cone having for its axis the tangent to the hyperbola; viz. the tangent in question is equally inclined to the two lines joining the point with the foci of the hyperbola, or with the extremities of the major axis of the ellipse. Taking $\theta$ for the inclination of the tangent to either of these lines, viz. $\theta$ is the semi-aperture of the cone, and $\gamma$ for the inclination of the tangent to the axis of $z$, then it is easy to show that
\[
\frac{g}{f} = \frac{\cos\gamma}{\sqrt{\cos^2\gamma - \sin^2\theta}},
\]
and we thence have
\[
V = \frac{2\pi c}{\cos\gamma}\sqrt{\cos^2\gamma - \sin^2\theta},
\]
viz. the ellipse is here considered as the section of a right cone of semi-aperture $\theta$, the perpendicular distance from the vertex being $=c$, and the inclination of this distance to the axis of the cone being $=\gamma$; and this being so, the potential is then expressed by the last preceding equation.

It will be observed that, when $\gamma = \frac{\pi}{2} - \theta$, the section becomes a parabola, and the potential is infinite; for any larger value of $\gamma$, the section is a hyperbola, and the formula ceases to be applicable.
\end{art}

\begin{art}
I originally obtained the result by thus considering the ellipse as the section of a right cone. Consider for a moment, in the case of any cone whatever, the plate included between the plane, perpendicular distance from the vertex $=c$, and the consecutive parallel plane, distance $=c+dc$. Let $dS$ denote an element of the first plane, $r$ its distance from the vertex, and $r + dr$ the distance produced to meet the second plane; also let $d\omega$ denote the subtended solid angle. We have
\[ dS \cdot dc = r^2\,dr\,d\omega, \]
or, since $\dfrac{dc}{dr} = \dfrac{c}{r}$, we obtain $dS = \dfrac{r^3}{c}\,d\omega$, or $-dS = -\dfrac{r^3}{c}\,d\omega$; wherefore the potential of the plane section is
\[
V = \frac{1}{c}\iint r\,d\omega,
\]
where $r$ denotes the value at a point of the plane section, and the integration extends over the spherical aperture of the cone.
\end{art}

\begin{art}
Let the position of $r$ be determined by means of its inclination $\theta$ to the axis of the cone, and the azimuth $\phi$ of the plane through $r$ and the axis of the cone; viz. taking the axis of the cone for the axis of $z$, suppose, as usual, $x = r\sin\theta\cos\phi$, $y = r\sin\theta\sin\phi$, $z = r\cos\theta$. We have then, as usual, $d\omega = \sin\theta\,d\theta\,d\phi$; and if the equation of the plane be
\[ x\cos\alpha + y\cos\beta + z\cos\gamma = c, \]
then the value of $r$ is obtained from the equation
\[ r\{(\cos\alpha\cos\phi + \cos\beta\sin\phi)\sin\theta + \cos\gamma\cos\theta\} = c; \]
so that we have for the potential
\[
V = c\iint \frac{\sin\theta\,d\theta\,d\phi}{\{(\cos\alpha\cos\phi + \cos\beta\sin\phi)\sin\theta + \cos\gamma\cos\theta\}^2},
\]
where the integration in regard to $\theta$ is from $\theta = 0$ to $\theta = $ the semi-aperture of the cone.
\end{art}

\begin{art}
The integral in regard to $\theta$ is
\[
\int \frac{\sin\theta\,d\theta}{(A\sin\theta + B\cos\theta)^2},
\]
where $A = \cos\alpha\cos\phi + \cos\beta\sin\phi$, $B = \cos\gamma$. Writing $A = R\sin N$, $B = R\cos N$, so that $R = \sqrt{A^2 + B^2}$, $\tan N = \dfrac{A}{B}$, the integral is
\[
\frac{1}{R^2}\int \frac{\sin\theta\,d\theta}{\sin^2(\theta + N)} = \frac{1}{R^2}\int \frac{\sin(\psi - N)\,d\psi}{\sin^2\psi},
\]
where $\psi = \theta + N$, $= \dfrac{1}{R^2}(-\cos N\cot\psi - \sin N\log\sin\psi)$; whence, between the limits $\theta = 0$ and $\theta = \Theta$, it is
\begin{multline*}
= \frac{1}{R^2}\cos N\{\cot N - \cot(\Theta + N)\} \\
+ \frac{1}{R^2}\sin N\{\log\sin(\Theta + N) - \log\sin N\}.
\end{multline*}
\end{art}

\begin{art}
But for the right cone, $\alpha = \beta = \frac{\pi}{2}$, $\cos\alpha = \cos\beta = 0$, $\cos\gamma = 1$, so that $A = 0$, $R = 1$, $\sin N = 0$, $\cos N = 1$, and the expression becomes simply $\cot N - \cot(\Theta + N)$, which between the limits is $= \cot\Theta - \cot\Theta = 0$, and the integral vanishes. This agrees with the known result that the potential of a circular plate on a point in its axis is simply proportional to the solid angle.
\end{art}

\begin{art}
Returning to the case of the right cone, we have $\cos\alpha = \cos\beta = 0$, $\cos\gamma = 1$, $R = 1$, $N = 0$, and the formula becomes
\[
V = c\int_0^{2\pi} d\phi\,\frac{\sin\Theta}{\cos\Theta + \sin\Theta\cos\phi},
\]
which can be at once evaluated. Writing $\cos\phi = \dfrac{1 - t^2}{1 + t^2}$, $\sin\phi = \dfrac{2t}{1 + t^2}$, $d\phi = \dfrac{2\,dt}{1 + t^2}$, the integral becomes
\[
V = 2c\int_{-\infty}^{\infty} \frac{\sin\Theta\,dt}{(1 + t^2)\cos\Theta + 2t\sin\Theta},
\]
which is readily found to be $= 2\pi c(1 - \cos\Theta)$, or restoring $\Theta = $ semi-aperture, we have the potential $= 2\pi c(1 - \cos\Theta)$, where $c$ is the distance of the vertex from the plane.
\end{art}

\begin{art}
For the ellipse considered as the section of a right cone of semi-aperture $\Theta$, the perpendicular distance from the vertex being $=c$, and the inclination of this distance to the axis of the cone being $=\gamma$, we have
\[
\frac{g}{f} = \frac{\cos\gamma}{\sqrt{\cos^2\gamma - \sin^2\Theta}},
\quad
\sin\Theta = \sqrt{\cos^2\gamma - \frac{g^2}{f^2}};
\]
and then, in place of the original formula, we may write
\[
V = \frac{2\pi c}{\cos\gamma}\sqrt{\cos^2\gamma - \sin^2\Theta},
\]
as already obtained.
\end{art}


\begin{art}
We may also obtain the result by a direct evaluation of the integral. The potential of an elliptic plate on an external point being
\[
V = 2fg\int_{\theta}^{\infty} \frac{dt}{\sqrt{t(t+f^2)(t+g^2)}},
\]
where $\theta$ is the positive root of
\[
\frac{a^2}{\theta+f^2} + \frac{b^2}{\theta+g^2} + \frac{c^2}{\theta} = 1,
\]
let the attracted point be on the focal hyperbola, so that $b=0$, $\dfrac{a^2}{f^2-g^2}=1$; then $\theta = \dfrac{f^2c^2}{g^2}$, and
\[
V = 2fg\int_{\frac{f^2c^2}{g^2}}}^{\infty} \frac{dt}{\sqrt{t(t+f^2)(t+g^2)}}.
\]
Writing $t = f^2\tan^2\phi$, this becomes
\begin{align*}
V &= 4fg\int_{\tan^{-1}\frac{c}{g}}^{\frac{\pi}{2}} \frac{d\phi}{\sqrt{f^2\sin^2\phi + g^2\cos^2\phi}} \\
&= 4g\int_{\tan^{-1}\frac{c}{g}}^{\frac{\pi}{2}} \frac{d\phi}{\sqrt{\sin^2\phi + \frac{g^2}{f^2}\cos^2\phi}}.
\end{align*}
But when $f$ varies and the point remains on the focal hyperbola at altitude $c$, we have $a^2 = f^2 - g^2$, and the value $V$ is independent of $f$; taking $f = \infty$, we find
\[
V = 4g\int_{\tan^{-1}\frac{c}{g}}^{\frac{\pi}{2}} \frac{d\phi}{\sin\phi}
= 4g\left[-\log\cot\frac{\phi}{2}\right]_{\tan^{-1}\frac{c}{g}}^{\frac{\pi}{2}}
= 2g\log\frac{\sqrt{c^2+g^2}+g}{\sqrt{c^2+g^2}-g},
\]
or as before $V = 2\pi(\sqrt{c^2+g^2} - c)$.
\end{art}

\begin{art}
The potential of the circle may be obtained directly. Taking the circle $x^2 + y^2 = g^2$, the potential on the point $(0, 0, c)$ is
\[
V = \iint \frac{dx\,dy}{\sqrt{x^2 + y^2 + c^2}}
= \int_0^{2\pi} d\phi \int_0^g \frac{r\,dr}{\sqrt{r^2 + c^2}}
= 2\pi(\sqrt{g^2 + c^2} - c),
\]
as already found.
\end{art}

\begin{art}
The theorem may be stated as follows: If a series of confocal ellipses be described, and from any point of the focal hyperbola as vertex a right cone be drawn having the ellipse for its section by a plane perpendicular to the tangent at the vertex; then the solid angle of the cone, multiplied by the distance of the vertex from the plane, is constant (in fact $= 2\pi(\sqrt{c^2+g^2} - c)$).
\end{art}

\begin{art}
We may verify this by direct calculation of the solid angle. The element of solid angle is $d\omega = \sin\theta\,d\theta\,d\phi$, and the whole solid angle is
\[
\omega = \int_0^{2\pi} d\phi \int_0^{\Theta} \sin\theta\,d\theta
= 2\pi(1 - \cos\Theta),
\]
where $\Theta$ is the semi-aperture. The potential is then $V = c\omega = 2\pi c(1 - \cos\Theta)$, which agrees with the former result.
\end{art}

\begin{art}
The potential of the ellipse may also be expressed in terms of elliptic functions. Writing for shortness
\[
k^2 = \frac{f^2 - g^2}{f^2},
\quad
\sin^2\psi = \frac{f^2}{\theta + f^2},
\]
then
\[
V = \frac{4fg}{f}\int_0^{\psi} \frac{d\psi}{\sqrt{1 - k^2\sin^2\psi}}
= 4gF(k, \psi),
\]
where $F(k, \psi)$ is the elliptic integral of the first kind.
\end{art}

\begin{art}
In the case where the attracted point is in the plane of the ellipse, $c=0$, we have $\theta = 0$ if $a^2 + b^2 < f^2$ (interior point), or $\theta$ positive if $a^2 + b^2 > f^2$ (exterior point). For an interior point, the potential is
\[
V = 2fg\int_0^{\infty} \frac{dt}{\sqrt{t(t+f^2)(t+g^2)}}
= 4gK(k),
\]
where $K(k)$ is the complete elliptic integral of the first kind.
\end{art}

\begin{art}
For a point on the axis of the ellipse, $a=0$, $b=0$, the equation for $\theta$ becomes
\[
\frac{c^2}{\theta} = 1,
\quad \text{or} \quad \theta = c^2,
\]
and the potential is
\[
V = 2fg\int_{c^2}^{\infty} \frac{dt}{\sqrt{t(t+f^2)(t+g^2)}}.
\]
If $f=g$, the ellipse becomes a circle, and we find
\[
V = 2\pi(\sqrt{c^2 + f^2} - c),
\]
as before.
\end{art}

\begin{art}
For the general case, we may express the result as follows. Let $e$ be the eccentricity of the ellipse, so that $g^2 = f^2(1-e^2)$; then
\[
V = 4f(1-e^2)^{\frac{1}{2}}\int_0^{\psi} \frac{d\psi}{\sqrt{1 - e^2\sin^2\psi}},
\]
where $\sin^2\psi = \dfrac{f^2}{\theta + f^2}$, and $\theta$ is determined by
\[
\frac{a^2}{\theta + f^2} + \frac{b^2}{\theta + f^2(1-e^2)} + \frac{c^2}{\theta} = 1.
\]
\end{art}

\begin{art}
In particular, if the attracted point is on the normal at the end of the major axis, $a=f$, $b=0$, then
\[
\frac{f^2}{\theta + f^2} + \frac{c^2}{\theta} = 1,
\]
giving $\theta^2 + f^2\theta - c^2f^2 - c^2\theta = \theta^2 + (f^2-c^2)\theta - c^2f^2 = 0$, or
\[
\theta = \frac{-(f^2-c^2) + \sqrt{(f^2-c^2)^2 + 4c^2f^2}}{2}
= \frac{-(f^2-c^2) + \sqrt{(f^2+c^2)^2}}{2}
= \frac{-(f^2-c^2) + (f^2+c^2)}{2} = c^2,
\]
or $\theta = -f^2$ which is inadmissible; so $\theta = c^2$, and
\[
\sin^2\psi = \frac{f^2}{c^2 + f^2},
\quad
\cos^2\psi = \frac{c^2}{c^2 + f^2}.
\]
\end{art}

\begin{art}
Suppose, to fix the ideas, $f=1$, and consider the points $(0, c)$ and $(a, 0)$, which have equal potentials. First, if $a > f$ (that is $a > 1$), the point $(a, 0)$ is outside the ellipse, and its potential is less than that of the centre; while the point $(0, c)$, if $c$ be properly chosen, may have any potential from $0$ to $\infty$. The condition for equal potentials gives a relation between $a$ and $c$.
\end{art}

\begin{art}
For the point $(0, c)$, on the minor axis, the equation for $\theta$ is
\[
\frac{c^2}{\theta} + \frac{0}{\theta + f^2} + \frac{0}{\theta + g^2} = 1,
\quad \text{i.e.} \quad \theta = c^2.
\]
The potential is
\[
V_1 = 2fg\int_{c^2}^{\infty} \frac{dt}{\sqrt{t(t+f^2)(t+g^2)}}.
\]
For the point $(a, 0)$, on the major axis, the equation is
\[
\frac{a^2}{\theta + f^2} = 1,
\quad \theta = a^2 - f^2.
\]
The potential is
\[
V_2 = 2fg\int_{a^2-f^2}^{\infty} \frac{dt}{\sqrt{t(t+f^2)(t+g^2)}}.
\]
For these to be equal, we must have $c^2 = a^2 - f^2$, or $c = \sqrt{a^2 - f^2}$.
\end{art}

\begin{art}
This result may be verified directly. If $c = \sqrt{a^2 - f^2}$, then for the point $(0, c)$:
\[
V_1 = 2fg\int_{a^2-f^2}^{\infty} \frac{dt}{\sqrt{t(t+f^2)(t+g^2)}},
\]
and for the point $(a, 0)$:
\[
V_2 = 2fg\int_{a^2-f^2}^{\infty} \frac{dt}{\sqrt{t(t+f^2)(t+g^2)}},
\]
so that $V_1 = V_2$, as stated.
\end{art}

\begin{art}
The potential of the ellipse may be expressed in a different form by means of the formula
\[
\int_{\theta}^{\infty} \frac{dt}{\sqrt{t(t+f^2)(t+g^2)}}
= \frac{2}{\sqrt{\theta + f^2}}F\left(\sqrt{\frac{f^2-g^2}{\theta+f^2}}, \frac{\pi}{2}\right),
\]
when $\theta > 0$. For an external point, the potential is thus expressed in terms of the complete elliptic integral of the first kind.
\end{art}

\begin{art}
The results obtained for the ellipse may be extended to the ellipsoid. The potential of an ellipsoid with semiaxes $(f, g, h)$ on the external point $(a, b, c)$ is
\[
V = 2fgh\int_{\theta}^{\infty} \frac{dt}{\sqrt{(t+f^2)(t+g^2)(t+h^2)}},
\]
where $\theta$ is the positive root of
\[
\frac{a^2}{\theta+f^2} + \frac{b^2}{\theta+g^2} + \frac{c^2}{\theta+h^2} = 1.
\]
This is the well-known formula of Ivory and Maclaurin.
\end{art}

\begin{art}
For the circle, putting $f=g$ in the formula for the ellipse, the equation for $\theta$ becomes
\[
\frac{a^2+b^2}{\theta+f^2} + \frac{c^2}{\theta} = 1,
\]
and the potential is
\[
V = 2f^2\int_{\theta}^{\infty} \frac{dt}{(t+f^2)\sqrt{t}}
= 4f^2\left[\frac{1}{f}\tan^{-1}\frac{\sqrt{t}}{f}\right]_{\theta}^{\infty}
= 2\pi f - 4f\tan^{-1}\frac{\sqrt{\theta}}{f}.
\]
\end{art}

\begin{art}
For the axial point $(0, 0, c)$, we have $\theta = c^2$ if we write $c$ for the distance, and then
\[
V = 2\pi f - 4f\tan^{-1}\frac{c}{f}
= 4f\left(\frac{\pi}{2} - \tan^{-1}\frac{c}{f}\right)
= 4f\tan^{-1}\frac{f}{c},
\]
or equivalently $V = 2\pi(\sqrt{c^2+f^2} - c)$, which is the same as before.
\end{art}

\begin{art}
The case of the circle may be treated independently by direct integration. Taking the circle $x^2 + y^2 = f^2$, and the attracted point $(0, 0, c)$, the potential is
\[
V = \int_0^{2\pi} d\phi \int_0^f \frac{r\,dr}{\sqrt{r^2 + c^2}}
= 2\pi\left[\sqrt{r^2+c^2}\right]_0^f
= 2\pi(\sqrt{f^2+c^2} - c).
\]
For an eccentric point $(a, b, c)$, with $a^2 + b^2 = d^2$, the potential may be shown to be
\[
V = 2\pi\sqrt{f^2 + c^2 + d^2 - 2fd\cos\psi} - 2\pi c,
\]
where $\psi$ is an auxiliary angle; or, more simply,
\[
V = 2\pi(\sqrt{R^2 + c^2} - c),
\]
where $R$ is the greatest distance of the point from the circumference of the circle.
\end{art}


% ====================================================================
% PAPER 604
% ====================================================================
\newpage
\section[Determination of the Attraction of an Ellipsoidal Shell on an Exterior Point]%
{604. Determination of the Attraction of an Ellipsoidal Shell on an Exterior Point.}

\begin{center}
{\small [From the \textit{Proceedings of the London Mathematical Society}, vol.~VI. (1874--1875), pp.~58--67.\\
Read January 14, 1875.]}
\end{center}

\begin{art}
The attraction of an ellipsoidal shell upon an exterior point may be determined by a method similar to that employed for the solid ellipsoid. Consider a homogeneous shell bounded by two similar and similarly situated concentric ellipsoids; say the outer ellipsoid has semiaxes $(f, g, h)$, and the inner ellipsoid has semiaxes $(f', g', h')$, where $\dfrac{f'}{f} = \dfrac{g'}{g} = \dfrac{h'}{h} = m$; the thickness of the shell is variable, being proportional to the distance from the centre.
\end{art}

\begin{art}
Let $P = (a, b, c)$ be an exterior point, and draw through $P$ a line cutting the ellipsoid; let this line meet the outer ellipsoid in $R'$ and $R''$, and the inner ellipsoid in $r'$ and $r''$. Let $Q$ be the point where the line meets the polar plane of $P$ with respect to the outer ellipsoid. Then the attractions of the two ellipsoids upon $P$ are inversely proportional to the squares of their masses, and the attraction of the shell is the difference of these attractions.
\end{art}

\begin{art}
Taking a confocal ellipsoid through $P$, let its semiaxes be $(F, G, H)$; then the potential of the outer ellipsoid at $P$ is proportional to $\dfrac{1}{F G H}$, and the attraction is directed towards the centre. The shell may be regarded as composed of elemental shells, each bounded by similar ellipsoids, and the total attraction is the integral of the attractions of these elemental shells.
\end{art}

\begin{art}
Let the equation of the outer ellipsoid be
\[
\frac{x^2}{f^2} + \frac{y^2}{g^2} + \frac{z^2}{h^2} = 1,
\]
and let the inner ellipsoid be
\[
\frac{x^2}{f'^2} + \frac{y^2}{g'^2} + \frac{z^2}{h'^2} = 1,
\]
where $f' = mf$, $g' = mg$, $h' = mh$, with $m < 1$. The shell is the region between these two surfaces.
\end{art}

\begin{art}
Consider a ray from $P$ meeting the outer ellipsoid at $R'$ and $R''$, and the inner at $r'$ and $r''$. Let $\rho$ be the distance from $P$ to a point on the ray, and let $\delta\rho$ be an element of length. The density being uniform, the mass of the element of the shell contained between $\rho$ and $\rho + d\rho$ is proportional to the area of the cross-section; and the attraction of this element on $P$ is inversely proportional to $\rho^2$.
\end{art}

\begin{art}
Let $\lambda, \mu, \nu$ be the direction-cosines of the ray from $P$; then a point on the ray is $(a + \lambda\rho, b + \mu\rho, c + \nu\rho)$. This point lies on the outer ellipsoid if
\[
\frac{(a + \lambda\rho)^2}{f^2} + \frac{(b + \mu\rho)^2}{g^2} + \frac{(c + \nu\rho)^2}{h^2} = 1,
\]
or, writing
\[
A = \frac{\lambda^2}{f^2} + \frac{\mu^2}{g^2} + \frac{\nu^2}{h^2},\quad
B = \frac{a\lambda}{f^2} + \frac{b\mu}{g^2} + \frac{c\nu}{h^2},\quad
C = \frac{a^2}{f^2} + \frac{b^2}{g^2} + \frac{c^2}{h^2} - 1,
\]
we have $A\rho^2 + 2B\rho + C = 0$, giving the distances $\rho_1$, $\rho_2$ to the points $R'$, $R''$.
\end{art}

\begin{art}
Similarly, for the inner ellipsoid, writing
\[
A' = \frac{\lambda^2}{f'^2} + \frac{\mu^2}{g'^2} + \frac{\nu^2}{h'^2}
= \frac{1}{m^2}A,\quad
B' = \frac{1}{m^2}B,\quad
C' = \frac{1}{m^2}C + \frac{1}{m^2} - 1,
\]
the distances to $r'$, $r''$ are given by $A'\rho^2 + 2B'\rho + C' = 0$.
\end{art}

\begin{art}
The attraction of the shell on $P$ in the direction $(\lambda, \mu, \nu)$ is found by integrating the attraction of the elemental slice. If $d\omega$ is the element of solid angle, the resolved attraction is
\[
dX = \lambda\,d\omega\int \rho\,d\rho,
\]
the integral being taken over the material of the shell along the ray. This gives
\[
dX = \lambda\,d\omega \cdot \frac{1}{2}(R'^2 - R''^2 - r'^2 + r''^2),
\]
where $R'$, $R''$ are the roots for the outer surface, and $r'$, $r''$ for the inner surface.
\end{art}

\begin{art}
The mass of the shell being
\[
M = \frac{4\pi}{3}fgh(1 - m^3),
\]
the attraction of the shell is found to be
\[
X = \frac{3M}{4\pi fgh(1-m^3)} \iint \lambda\,(R'^2 - R''^2 - r'^2 + r''^2)\,d\omega.
\]
By known properties of the ellipsoid, this reduces to a simple form involving only elementary functions.
\end{art}

\begin{art}
The final result for the attraction of the shell on the exterior point $P$ is as follows. Let $\theta$ be the positive root of
\[
\frac{a^2}{\theta + f^2} + \frac{b^2}{\theta + g^2} + \frac{c^2}{\theta + h^2} = 1,
\]
and let $(F, G, H)$ be the semiaxes of the confocal ellipsoid through $P$, so that $F^2 = \theta + f^2$, $G^2 = \theta + g^2$, $H^2 = \theta + h^2$. Then the attraction of the shell is the same as that of a solid ellipsoid of semiaxes $(F, G, H)$ and density adjusted so that the mass is equal to that of the shell; viz. the attraction is directed towards the centre, and its magnitude is
\[
\frac{3M}{\sqrt{(\theta+f^2)(\theta+g^2)(\theta+h^2)}} \cdot \frac{1}{\theta^{\frac{3}{2}}}
\times \text{(proper numerical factor)}.
\]
\end{art}

\begin{art}
Analytical Expressions for the Attraction of the Shell, and for the Resolved Attractions.
\end{art}

\begin{art}
The attraction of the shell was shown to be expressible in terms of the confocal ellipsoid through the attracted point. If the semiaxes of this confocal be $(F, G, H)$, then the attraction of the shell is proportional to the attraction of the solid ellipsoid with semiaxes $(F, G, H)$; and the resolved attractions are
\[
X = -\frac{3Ma}{\theta}\cdot\frac{1}{F G H},\quad
Y = -\frac{3Mb}{\theta}\cdot\frac{1}{F G H},\quad
Z = -\frac{3Mc}{\theta}\cdot\frac{1}{F G H},
\]
where $\theta$ is determined as above.
\end{art}

\begin{art}
The mass of the shell is $\dfrac{4\pi}{3}fgh \cdot S_m$, where $S_m$ is a factor depending on the thickness; whence
\[
\text{Resolved Attraction} = \frac{M}{\text{Mass of confocal}} \times \text{Attraction of confocal},
\]
and since the attraction of the confocal solid ellipsoid is known, we have the result.
\end{art}

\begin{art}
The direction-cosines of the attraction are proportional to $\dfrac{a}{\theta+f^2}$, $\dfrac{b}{\theta+g^2}$, $\dfrac{c}{\theta+h^2}$; and the resultant attraction is
\[
R = \sqrt{X^2 + Y^2 + Z^2}
= \frac{3M}{F G H}\sqrt{\frac{a^2}{(\theta+f^2)^2} + \frac{b^2}{(\theta+g^2)^2} + \frac{c^2}{(\theta+h^2)^2}}.
\]
\end{art}

\begin{art}
For a spherical shell, $f=g=h$, the formula simplifies considerably. The confocal through $P$ is a concentric sphere of radius $\sqrt{\theta + f^2}$, where $\theta = a^2 + b^2 + c^2 - f^2 = d^2 - f^2$, $d$ being the distance of $P$ from the centre. The attraction is then
\[
X = -\frac{M}{d^2}\cdot\frac{a}{d},\quad
Y = -\frac{M}{d^2}\cdot\frac{b}{d},\quad
Z = -\frac{M}{d^2}\cdot\frac{c}{d},
\]
which is Newton's well-known result: the spherical shell attracts as if its mass were concentrated at the centre.
\end{art}


% ====================================================================
% PAPER 605
% ====================================================================
\newpage
\section[Note on a Point in the Theory of Attraction]%
{605. Note on a Point in the Theory of Attraction.}

\begin{center}
{\small [From the \textit{Proceedings of the London Mathematical Society}, vol.~VI. (1874--1875), pp.~79--81.\\
Read February 11, 1875.]}
\end{center}

\begin{art}
Consider a mass of matter distributed in any manner throughout space, and let $V$ be the potential at any point. The potential satisfies Laplace's equation $\nabla^2 V = 0$ at all points external to the attracting mass, and Poisson's equation $\nabla^2 V = -4\pi\rho$ at points where the density is $\rho$.
\end{art}

\begin{art}
It is important to observe that the potential and its first derivatives are everywhere continuous, even at the surface of the attracting body. The second derivatives, however, are discontinuous at the surface; in fact, if $V_1$ is the potential outside and $V_2$ inside, then at the surface
\[
V_1 = V_2, \quad \frac{\partial V_1}{\partial n} = \frac{\partial V_2}{\partial n},
\]
but
\[
\frac{\partial^2 V_1}{\partial n^2} - \frac{\partial^2 V_2}{\partial n^2} = -4\pi\sigma,
\]
where $\sigma$ is the surface-density.
\end{art}

\begin{art}
For a spherical shell; the potential is a different function for external and interior points, viz. for internal points it is a constant, $= M \div$ radius; for external points it is $= M \div$ distance from the centre. But at the surface itself, the two expressions coincide, and the first derivative is continuous. The second derivative, however, has a discontinuity proportional to the surface-density.
\end{art}

\begin{art}
This is best shown by the annexed figure, which represents a uniform spherical shell made up of two segments, one of which is taken to be a small segment at the point of observation. The potential due to the small segment varies continuously, but its normal derivative changes abruptly by $4\pi\sigma$ on crossing the surface.
\end{art}

\begin{art}
The general theorem may be stated as follows: If $V$ be the potential of any distribution of matter, and if we cross a surface at which there is a stratum of matter of surface-density $\sigma$, then the normal derivative of $V$ changes by $-4\pi\sigma$; that is,
\[
\left(\frac{\partial V}{\partial n}\right)_1 - \left(\frac{\partial V}{\partial n}\right)_2 = -4\pi\sigma,
\]
where the suffixes $1$ and $2$ refer to the two sides of the surface. This is the well-known theorem of Green.
\end{art}

\begin{art}
The proof is immediate from the consideration of the elemental tube of force. Taking a small cylinder with its axis along the normal to the surface, and its two ends on opposite sides of the surface, Gauss's theorem gives at once the discontinuity of the normal force. Since the potential itself is continuous, and since the tangential derivatives are also continuous (being determined solely by the values of $V$ on the surface), the theorem is completely established.
\end{art}


% ====================================================================
% PAPER 606
% ====================================================================
\newpage
\section[On the Expression of the Coordinates of a Point of a Quartic Curve as Functions of a Parameter]%
{606. On the Expression of the Coordinates of a Point of a Quartic Curve as Functions of a Parameter.}

\begin{center}
{\small [From the \textit{Proceedings of the London Mathematical Society}, vol.~VI. (1874--1875), pp.~81--83.\\
Read February 11, 1875.]}
\end{center}

\begin{art}
The question of expressing the coordinates of a point on a plane curve as rational functions of a parameter is one of great importance in the theory of curves. For a conic, the solution is well known; for a cubic, it can be effected by means of elliptic functions. For a quartic curve, the problem is more complicated, but can be solved in certain cases by means of elliptic functions.
\end{art}

\begin{art}
Consider a quartic curve having a double point; such a curve may be regarded as the reciprocal of a cubic. Let the equation of the quartic be
\[
U = (ax + by + c)^2(x^2 + y^2 + 1) + 2(ax + by + c)(dx^2 + exy + fy^2 + gx + hy + k) + \cdots = 0.
\]
By a suitable transformation, this may be reduced to a standard form.
\end{art}

\begin{art}
Let the quartic have a node at the origin; then the equation may be written
\[
(ax^2 + 2hxy + by^2)^2 + 2(ax + by + c)(a'x^2 + 2h'xy + b'y^2) + \cdots = 0.
\]
Taking a line $y = tx$ through the node, the other intersections are given by a quadratic equation, whose roots determine two points on the curve. If these coincide, the line is a tangent; but in general, we obtain two points corresponding to each value of $t$.
\end{art}

\begin{art}
The condition for the line $y = tx$ to meet the quartic in two coincident points (other than the node) gives a quartic equation in $t$; but by introducing an auxiliary parameter, the roots of this quartic may be expressed in terms of elliptic functions. In fact, the quartic equation in $t$ may be written
\[
At^4 + Bt^3 + Ct^2 + Dt + E = 0,
\]
and by the substitution $t = \tan\phi$, this may be reduced to the form
\[
P\cos 2\phi + Q\sin 2\phi + R\cos 4\phi + S\sin 4\phi = T,
\]
which is solvable by elliptic functions.
\end{art}

\begin{art}
Writing for a moment $ax + b = 2\sqrt{b^2-ac} \cdot u$, the equation becomes
\[
4u^2 - 3u + \frac{a^2d - 3abc + 2b^3}{2(b^2-ac)\sqrt{b^2-ac}} = 0.
\]
Hence, assuming
\[
\frac{a^2d - 3abc + 2b^3}{2(b^2-ac)\sqrt{b^2-ac}} = \cos\phi,
\]
we have
\[
4u^2 - 3u + \cos\phi = 0,
\]
which gives $u$ in terms of $\phi$.
\end{art}

\begin{art}
The roots of this cubic in $u$ are connected with the trisection of the angle; in fact, $u = \cos\frac{\phi + 2k\pi}{3}$ for $k = 0, 1, 2$. The three values of $u$ correspond to the three real intersections of the line with the cubic curve.
\end{art}

\begin{art}
For the quartic curve, we proceed as follows. Let the equation be
\[
y^2 = (x - \alpha)(x - \beta)(x - \gamma)(x - \delta),
\]
which is the standard form of a quartic with all roots real. By the substitution
\[
x = \frac{\alpha(x' - \gamma) - \gamma(x' - \alpha)}{(x' - \gamma) - (x' - \alpha)}
= \frac{(\alpha - \gamma)x'}{\alpha - \gamma},
\]
this may be reduced to the form
\[
y'^2 = x'(1 - x')(1 - k^2x'),
\]
which is the elliptic normal form. The coordinates are then expressible as elliptic functions of a parameter.
\end{art}

\begin{art}
In general, for any plane quartic curve, the coordinates may be expressed as elliptic functions of a parameter by means of the following construction. Take two fixed points $A$, $B$ on the curve, and let $P$ be a variable point. The line $AP$ meets the curve again in $P'$, and the line $BP'$ meets again in $P''$, and so on. The correspondence between $P$ and $P''$ is a $(1,1)$ correspondence, and the parameter of $P''$ differs from that of $P$ by a constant. Hence the parameter of $P$ may be taken as an elliptic parameter.
\end{art}

\begin{art}
The result may be stated more precisely: A plane quartic curve is of deficiency $3$; but if it have a double point, the deficiency is reduced to $2$, and if it have two double points, to $1$. In the latter case, the curve is unicursal, and the coordinates may be expressed rationally in terms of a parameter. If the deficiency is $2$, the coordinates may be expressed in terms of elliptic functions. For the general quartic (deficiency $3$), the coordinates require hyperelliptic functions of two arguments.
\end{art}

\begin{art}
And consequently, for the case of a single double point (deficiency $2$), the coordinates $(x, y)$ may be expressed as elliptic functions of a single parameter $u$. The relation between $x$ and $y$ is given by the equation of the curve, and either coordinate may be taken as the independent elliptic variable. The periods of the elliptic functions are determined by the moduli, which are algebraic functions of the coefficients of the quartic.
\end{art}


% ====================================================================
% PAPER 607
% ====================================================================
\newpage
\section[A Memoir on Prepotentials]%
{607. A Memoir on Prepotentials.}

\begin{center}
{\small [From the \textit{Philosophical Transactions of the Royal Society of London}, vol.~CLXV. Part II. (1875), pp.~675--774.\\
Received April 8,---Read June 10, 1875.]}
\end{center}

\subsection*{Introduction.}

\begin{art}
The present memoir is concerned with the theory of ``prepotentials,'' a term which I use to denote a class of integrals generalizing the ordinary potential of attracting matter. The potential of a body on an external point depends on the inverse distance; the prepotential involves an arbitrary index $q$, and reduces to the potential for particular values of $q$.
\end{art}

\begin{art}
The ordinary potential of a body of density $\rho$ at a point $(a, b, c)$ is
\[
V = \iiint \frac{\rho\,dx\,dy\,dz}{\sqrt{(a-x)^2 + (b-y)^2 + (c-z)^2}}.
\]
More generally, the prepotential is defined as
\[
W = \iiint \frac{\rho\,dx\,dy\,dz}{\{(a-x)^2 + (b-y)^2 + (c-z)^2\}^{\frac{1}{2}q}},
\]
where $q$ is any positive or negative number. When $q=1$, this reduces to the ordinary potential; when $q=3$, it is the second polar; and so on.
\end{art}

\begin{art}
The prepotential may be considered in various cases according to the dimension of the attracting body:
\begin{enumerate}
\item[A.] The potential-point case; $q$ arbitrary, the attracting body being a single point.
\item[B.] The potential-curve case; $q$ arbitrary, the attracting body being a curve.
\item[C.] The potential-surface case; $q = -\frac{1}{2}$, the surface arbitrary, so that the integral is
\[
\iint \frac{\rho\,dS}{\{(a-x)^2 + \cdots + (c-z)^2 + (e-w)^2\}^{\frac{1}{2}}}.
\]
\item[D.] The potential-solid case; $q = -\frac{1}{2}$, the solid arbitrary.
\end{enumerate}
\end{art}

\begin{art}
The cases of a less-than-three-dimensional volume are in the present memoir considered only incidentally. It is sufficient to remark that the line-integral and the surface-integral may be deduced as limiting cases of the volume integral, by supposing the density to become infinite in a certain manner.
\end{art}

\begin{art}
The potential-surface integral, provided that the point $(a, \ldots, c, e)$ does not lie within the material space, satisfies Laplace's equation in the coordinates of the attracted point. I would rather say that the integral does not satisfy Laplace's equation when the attracted point is within the material space; for in this case the integral has a different form, depending on the side from which the surface is approached.
\end{art}

\subsection*{Case $q = $ Positive.}

\begin{art}
I consider first the case where $q$ is positive. The value is here
\[
W = \frac{1}{2\pi^2}\frac{\Gamma(\frac{1}{2}q)}{\Gamma(\frac{1}{2}n - \frac{1}{2}q)} \int_0^{\infty} \frac{u^{\frac{1}{2}n - \frac{1}{2}q - 1}\,du}{(1 + u)^{\frac{1}{2}n}},
\]
or, since $\dfrac{1}{2\pi^2}$ is indefinitely small when $n$ is large, we may write
\[
W = \frac{\Gamma(\frac{1}{2}q)}{\pi^{\frac{1}{2}n}} \int \frac{\rho\,d\omega}{r^q},
\]
where $d\omega$ is the element of volume, and $r$ the distance.
\end{art}

\begin{art}
The prepotential of a plane surface of uniform density on a point at distance $e$ from the plane is required in the case where the distance $e$ (taken to be always positive) is indefinitely small in regard to the radius $R$. Writing $x = r\xi$, $y = r\eta$, $z = r\zeta$, where $r$ is a length, the integral becomes
\[
W = \int_0^R \frac{r^{n-1}\,dr}{r^q} \int \frac{\rho\,d\Omega}{(1 + \text{terms in } \xi, \eta, \ldots)^{\frac{1}{2}q}},
\]
where $d\Omega$ is the element of solid angle.
\end{art}

\begin{art}
For a disk of radius $R$ and surface-density $\rho$, the prepotential at a point on the axis at distance $e$ is
\[
W = 2\pi\rho \int_0^R \frac{r\,dr}{(r^2 + e^2)^{\frac{1}{2}q}}.
\]
When $q=1$, this gives the ordinary potential; when $q > 1$, the integral converges even for $e=0$; and when $q < 1$, the value for $e=0$ requires special consideration.
\end{art}

\begin{art}
Evaluating the integral, we have
\[
W = \frac{2\pi\rho}{2-q}\left[\frac{1}{(R^2+e^2)^{\frac{1}{2}q-1}} - \frac{1}{e^{q-2}}\right]
\]
for $q \neq 2$; and for $q=2$,
\[
W = \pi\rho\log\frac{R^2+e^2}{e^2}.
\]
When $e$ is small compared with $R$, this becomes approximately
\[
W = \frac{2\pi\rho}{2-q}R^{2-q} \quad (q < 2),
\]
and for $q = 2$,
\[
W = 2\pi\rho\log\frac{R}{e}.
\]
\end{art}

\begin{art}
The prepotential of the remaining portion of the surface vanishes in comparison with that of the disk. But without entering into the question I assume that the prepotential of the whole surface may be obtained by integration, and that the result is finite for all values of $q$ less than the dimension of the space.
\end{art}

\begin{art}
The differential equation satisfied by the prepotential may be obtained as follows. We have
\[
\frac{\partial W}{\partial a} = q \iiint \frac{\rho(x-a)\,dx\,dy\,dz}{r^{q+2}},
\]
and
\[
\frac{\partial^2 W}{\partial a^2} = q \iiint \frac{\rho\{q(x-a)^2 - r^2\}\,dx\,dy\,dz}{r^{q+4}}.
\]
Adding the similar equations for $b$ and $c$, we obtain
\[
\nabla^2 W = q(q-1) \iiint \frac{\rho\,dx\,dy\,dz}{r^{q+2}},
\]
or
\[
\nabla^2 W = q(q-1) W_{q+2},
\]
where $W_{q+2}$ denotes the prepotential with index $q+2$.
\end{art}

\begin{art}
This equation connects the prepotentials of different orders. In particular, for the ordinary potential ($q=1$), we have $\nabla^2 W_1 = 0$ at external points; and for $q=3$, $\nabla^2 W_3 = 6 W_5$. These relations are useful in the solution of problems in attraction.
\end{art}

\begin{art}
The value of $\dfrac{\partial W}{\partial e}$, where $W$ is given as a function of $(x, \ldots, z, w)$, requires explanation. Taking $\cos\alpha$, $\ldots$, $\cos\gamma$ to be the inclinations of the normal to the coordinate axes, we have
\[
\frac{\partial W}{\partial e} = \cos\alpha \frac{\partial W}{\partial x} + \cdots + \cos\gamma \frac{\partial W}{\partial z},
\]
and this expression determines the rate of change of $W$ in the direction of the normal.
\end{art}

\begin{art}
The prepotential of a solid sphere of radius $R$ and uniform density $\rho$ on an external point at distance $f$ from the centre is
\[
W = \frac{4\pi\rho}{3-q}\frac{R^3}{f^q} \cdot {}_2F_1\left(\frac{1}{2}q, \frac{3}{2}; \frac{5}{2}; \frac{R^2}{f^2}\right),
\]
where ${}_2F_1$ is the hypergeometric function. For the ordinary potential ($q=1$), this reduces to $\dfrac{M}{f}$, where $M = \dfrac{4\pi\rho R^3}{3}$ is the mass.
\end{art}

\begin{art}
When the attracted point is inside the sphere, the result is different. We have
\[
W = \frac{4\pi\rho}{3-q}\left[\frac{R^{3-q}}{3-q} + \int_r^R \frac{r'^{2-q}\,dr'}{r'^{1-q}}\right],
\]
which for $q=1$ gives
\[
W = 2\pi\rho\left(R^2 - \frac{1}{3}r^2\right),
\]
the well-known formula for the internal potential of a uniform sphere.
\end{art}

\begin{art}
The general formula for the prepotential of a uniform solid sphere is
\[
W = \frac{4\pi\rho}{(3-q)(2-q)}\left[R^{2-q} - \frac{2-q}{2}r^2 R^{-q} + \cdots\right],
\]
where $\rho$ here denotes the density at the point $N$; and the value of the integral is
\[
W = \frac{1}{2}R^{2-q}(1 + \text{terms in } \xi, \eta, \ldots)^{-\frac{1}{2}q} \cdot \frac{\rho}{\text{etc.}}
\]
Observe that this is indefinitely small, as it should be, when the point is outside the sphere.
\end{art}

\begin{art}
The prepotential of a surface-distribution on a point near the surface may be treated similarly. Let $\sigma$ be the surface-density, and let the surface be plane near the attracted point. Then the prepotential is
\[
W = \sigma \iint \frac{dS}{r^q},
\]
and for a circular area of radius $R$,
\[
W = \frac{2\pi\sigma}{2-q}\left[(R^2+e^2)^{1-\frac{1}{2}q} - e^{2-q}\right].
\]
When $e = 0$, this becomes $W = \dfrac{2\pi\sigma R^{2-q}}{2-q}$ for $q < 2$; and for $q = 2$, $W = 2\pi\sigma\log R$.
\end{art}

\begin{art}
The discontinuity in the normal derivative of the prepotential at a surface may be found as follows. Let $W_+$ be the value on the positive side, and $W_-$ on the negative side; then
\[
\frac{\partial W_+}{\partial n} - \frac{\partial W_-}{\partial n} = -\frac{2\pi^{\frac{1}{2}n}\sigma}{\Gamma(\frac{1}{2}n - 1)} \cdot \frac{1}{e^{q-1}},
\]
where $(\cos\alpha', \ldots, \cos\gamma')$ and $(\cos\alpha'', \ldots, \cos\gamma'')$ are the cosine-inclinations of the normal on the two sides of the surface. For the ordinary potential ($q=1$), this reduces to the theorem of Green.
\end{art}

\begin{art}
Reverting to the original form of the prepotential, we have for the solid case
\[
W = \iiint \frac{\rho\,dx\,dy\,dz}{\{(a-x)^2 + (b-y)^2 + (c-z)^2\}^{\frac{1}{2}q}}.
\]
To get rid of the constant factors, we may write this as
\[
W = \frac{1}{\Gamma(\frac{1}{2}q)} \int_0^{\infty} \lambda^{\frac{1}{2}q-1}\,d\lambda \iiint \rho\,dx\,dy\,dz\,e^{-\lambda r^2},
\]
which is a useful form for theoretical investigations.
\end{art}

\begin{art}
By differentiation under the integral sign, we obtain formulae for the derivatives of the prepotential. For instance,
\[
\frac{\partial W}{\partial a} = \frac{q}{\Gamma(\frac{1}{2}q)} \int_0^{\infty} \lambda^{\frac{1}{2}q}\,d\lambda \iiint \rho(x-a)\,dx\,dy\,dz\,e^{-\lambda r^2},
\]
and similarly for the higher derivatives.
\end{art}

\begin{art}
The prepotential may be expanded in a series of spherical harmonics when the attracting body is a sphere or an ellipsoid. For a sphere of radius $R$, we have
\[
W = \frac{M}{f^q}\left[1 + \frac{q(q+2)}{6}\frac{R^2}{f^2} + \frac{q(q+2)(q+4)}{120}\frac{R^4}{f^4} + \cdots\right]
\]
for an external point at distance $f > R$; and for an internal point at distance $r < R$,
\[
W = \frac{M}{R^q}\left[1 + \frac{q(2-q)}{6}\frac{r^2}{R^2} + \frac{q(2-q)(4-q)}{120}\frac{r^4}{R^4} + \cdots\right].
\]
\end{art}

\begin{art}
The case $q = 0$ gives simply the mass of the body; $q = -1$ gives the moment of inertia about the attracted point; and higher negative values of $q$ give higher moments. The case $q = 2$ gives the potential of a line-distribution; $q = 3$ gives the potential of a surface-distribution; and so on.
\end{art}

\begin{art}
For the ellipsoid, with semiaxes $(f, g, h)$ and uniform density $\rho$, the prepotential on an external point $(a, b, c)$ is
\[
W = \pi\rho fgh \int_{\theta}^{\infty} \frac{dt}{\sqrt{(t+f^2)(t+g^2)(t+h^2)}} \cdot \frac{1}{(t + \theta)^{\frac{1}{2}q}},
\]
where $\theta$ is the positive root of the equation
\[
\frac{a^2}{\theta + f^2} + \frac{b^2}{\theta + g^2} + \frac{c^2}{\theta + h^2} = 1.
\]
This formula generalizes the well-known result of Ivory and Maclaurin.
\end{art}

\begin{art}
The prepotential has important applications in the theory of elasticity, in the theory of heat, and in electricity. In elasticity, the displacement due to a body-force may be expressed in terms of prepotentials; in the theory of heat, the temperature due to a variable source involves prepotential integrals; and in electricity, the potential of a variable charge-distribution is a prepotential.
\end{art}

\begin{art}
The memoir concludes with a table of the principal formulae for prepotentials in one, two, and three dimensions, for various values of the index $q$, and for the principal forms of attracting bodies: the point, the line, the disk, the sphere, and the ellipsoid. These formulae are useful in the solution of physical problems involving generalized potentials.
\end{art}


\end{document}
